\documentclass[a4paper,11pt]{article}

\usepackage[utf8]{inputenc}
\usepackage[portuguese]{babel}
\usepackage[a4paper,margin=2.5cm]{geometry}
\usepackage{graphicx}
\usepackage{xcolor}
\usepackage[colorlinks=true,linkcolor=black,urlcolor=blue,citecolor=black]{hyperref}
\usepackage{enumitem}
\usepackage{titlesec}
\usepackage{fancyhdr}

\definecolor{iselblue}{RGB}{0,51,102}

\titleformat{\section}{\large\bfseries\color{iselblue}}{\thesection}{0.6em}{}
\titlespacing*{\section}{0pt}{1.4em}{0.6em}

\pagestyle{fancy}
\fancyhf{}
\renewcommand{\headrulewidth}{0.6pt}
\fancyhead[L]{\small Instituto Superior de Engenharia de Lisboa}
\fancyhead[R]{\small EduCal --- Informações de Entrega}
\fancyfoot[C]{\small\thepage}

\begin{document}

\begin{center}
  \includegraphics[width=45mm]{./FinalReport/figures/logoISELnew2.png}\\[4mm]
  {\Large\bfseries\color{iselblue} EduCal --- Academic Planning Application}\\[2mm]
  {\large Informações de Entrega}\\[4mm]
  \begin{tabular}{rl}
    Autor: & Daniel Santos, nº 51701 \\
    Curso: & Licenciatura em Engenharia Informática, Redes e Telecomunicações \\
    Orientadores: & Filipe Freitas e Sandra Ferreira \\
    Data: & Julho de 2026 \\
  \end{tabular}
\end{center}

\vspace{6mm}
\noindent Este documento acompanha o Relatório Final (\texttt{FinalReport.pdf}) e tem como objetivo situar quem o recebe: onde está o código-fonte, onde pode ser experimentada a aplicação, e como o próprio relatório está organizado.

\section{Repositório de Código}
\url{https://github.com/Patilhas/educal}

\noindent O repositório contém a totalidade do código-fonte do projeto (frontend, backend e camada de dados), bem como o próprio código-fonte LaTeX deste relatório, em \texttt{docs/reports/FinalReport/}.

\section{Aplicação em Produção (Vercel)}
A aplicação encontra-se atualmente em execução numa demonstração pública, alojada na Vercel:

\begin{center}
  \url{https://educal-azure.vercel.app/}
\end{center}

\noindent \textbf{Utilizador de demonstração}:
\begin{center}
\begin{tabular}{rl}
  Email: & \texttt{daniel@educal.local} \\
  Password: & \texttt{Password123!} \\
\end{tabular}
\end{center}

\noindent \textit{Nota:} esta demonstração utiliza a implementação de dados assente em Redis (Upstash), conforme descrito na Secção 3.4 do relatório (``Data Access Strategy''), e não a base de dados PostgreSQL prevista para a instalação institucional final (Capítulo 6, ``Deployment and Infrastructure'').

\section{Estrutura do Relatório}
O relatório está organizado nos seguintes capítulos e apêndices:

\begin{description}[style=nextline, leftmargin=0pt]
  \item[Capítulo 1 --- Introduction] Contexto do projeto, objetivos e requisitos funcionais/não funcionais.
  \item[Capítulo 2 --- Background and Problem Definition] Descrição do processo atual (baseado em folhas de cálculo) e das suas limitações e estado da arte face a ferramentas de calendário genéricas.
  \item[Capítulo 3 --- System Architecture and Data Model] Arquitetura em três camadas, stack tecnológica (Next.js, PostgreSQL, Drizzle ORM, React), modelo de dados e diagrama Entidade-Relação, e a estratégia de acesso a dados (PostgreSQL/Redis, intercambiável).
  \item[Capítulo 4 --- Backend] Organização em camadas (Controller-Service-Data), autenticação e controlo de acesso baseado em papéis, tratamento de erros, o Motor de Migração, o Motor de Regras de Negócio, o sistema de notificações, e os testes automatizados.
  \item[Capítulo 5 --- Frontend] Organização do código do frontend, comunicação com o backend, e descrição de todos os ecrãs e diálogos da aplicação.
  \item[Capítulo 6 --- Deployment and Infrastructure] Estratégia de implementação prevista e infraestrutura atual (Vercel para demonstração, VM institucional do ISEL como destino final).
  \item[Capítulo 7 --- Future Work] Partes do âmbito inicial ainda não entregues (exportação para Excel, modelos de texto para email) e os próximos passos concretos para a instalação institucional.
  \item[Capítulo 8 --- Conclusion] Síntese dos requisitos entregues e balanço final do projeto.
  \item[Referências] Bibliografia citada ao longo do relatório.
  \item[Apêndice A --- Detailed Entity-Relationship Diagram] Versão completa do diagrama Entidade-Relação, com todas as colunas e respetivos tipos.
  \item[Apêndice B --- Selected Code Listings] Código-fonte integral de alguns componentes centrais, mostrados de forma abreviada nos capítulos principais.
  \item[Apêndice C --- Additional Screenshots] Capturas de ecrã adicionais das vistas do calendário e do modal de validação/migração, referidas no Capítulo 5.
\end{description}

\end{document}
